Use $A=(0,0)$, $B=(1,0)$, $C=(1,c)$, and $D=(d,-\dfrac{c}{4})$ with $c>0$ and $0<d<1$. Put $t=\tan\phi>0$. Then $t=\dfrac{c}{[4(1-d)]}$. The right angle at $D$ says
\[
(A-D)\cdot(C-D)=d^2-d+\dfrac{5c^2}{16}=0.
\]
Substituting $d=1-\dfrac{c}{(4t)}$ yields
\[
c=\dfrac{4t}{1+5t^2},\qquad d=\dfrac{5t^2}{1+5t^2}.
\]

In part 1, $c=\tan3\phi=\dfrac{(3t-t^3)}{(1-3t^2)}$, with $0<\phi<30^\circ$. Equating the expressions gives $5t^4-26t^2+1=0$. The root in $0<t^2<\dfrac{1}{3}$ is
\[
\boxed{\tan^2\phi=\dfrac{13-2\sqrt{41}}{5}}.
\]
In part 2, $c=\cot3\phi=\dfrac{(1-3t^2)}{(3t-t^3)}$. The resulting equation is $11t^4+10t^2-1=0$, whose positive root is
\[
\boxed{\tan^2\phi=\dfrac{1}{11}}.
\]
Both results give $0<d<1$ and positive $c$, and the angle bounds ensure the triple-angle equations select the required acute angles.
\begin{asy}
texpreamble("\usepackage{amsmath,amssymb}");
size(240);
defaultpen(fontsize(9pt));
pair A=(0,0),B=(1,0),C=(1,sqrt(11)/4),D=(5.0/16,-sqrt(11)/16),H=(5.0/16,0);draw(A--B--C--cycle);draw(A--D--C,blue);draw(B--D);draw(D--H,dashed);label("$A$",A,(-1,0));label("$B$",B,(1,0));label("$C$",C,(0,1));label("$D$",D,(0,-1));label("$H$",H,(0,1));
\end{asy}
